
\documentclass{article}
\usepackage{colortbl}
\usepackage{makecell}
\usepackage{multirow}
\usepackage{supertabular}

\begin{document}

\newcounter{utterance}

\centering \large Interaction Transcript for game `eqbench', experiment `eq\_bench', episode 13 with BSU{-}SLIM\_\_Qwen3.5{-}9B{-}Base\_\_playpen{-}prm{-}9b{-}best\_of\_n.
\vspace{24pt}

{ \footnotesize  \setcounter{utterance}{1}
\setlength{\tabcolsep}{0pt}
\begin{supertabular}{c@{$\;$}|p{.15\linewidth}@{}p{.15\linewidth}p{.15\linewidth}p{.15\linewidth}p{.15\linewidth}p{.15\linewidth}}
    \# & \multicolumn{2}{c}{Player} && \multicolumn{2}{c}{Game Master} \\
    \hline

    \theutterance \stepcounter{utterance}  
    & & & \multicolumn{4}{p{0.6\linewidth}}{
        \cellcolor[rgb]{0.9,0.9,0.9}{
            \makecell[{{p{\linewidth}}}]{
                \texttt{\tiny{[P1$\langle$GM]}}
                \\ 
\texttt{Your task is to predict the likely emotional responses of a character in this dialogue:} \\
\\ 
\texttt{Harold: I don't understand why you're so worked up about this election. It's not like the world is going to end if your candidate doesn't win.} \\
\texttt{Lorelei: That's not the point, Harold. It's about what this election represents. It's about the future of our planet, our society.} \\
\texttt{Harold: And you think your candidate has all the answers? } \\
\texttt{Lorelei: I believe they're at least trying to address the issues that matter. Climate change, social justice, healthcare...} \\
\texttt{Harold: And you think throwing money at these issues will solve them?} \\
\texttt{Lorelei: It's not about throwing money, Harold. It's about investing in solutions, in our future.} \\
\texttt{Harold: And who's going to pay for all these grand plans? Us, the taxpayers. } \\
\texttt{Lorelei: Isn't that the point of taxes? To fund public services and infrastructure?} \\
\texttt{Harold: Yes, but there's a limit. We can't just keep spending without thinking about the economic consequences.} \\
\texttt{Lorelei: And we can't keep ignoring the environmental and social consequences of our actions either.} \\
\texttt{Harold: I'm not suggesting we ignore them, Lorelei. I'm just saying we need to be realistic about what we can afford.} \\
\texttt{Lorelei: And I'm saying we can't afford not to act. The cost of inaction is far greater.} \\
\texttt{Harold: And who decides what action is necessary? The government? The same government that can't even balance a budget?} \\
\texttt{{[}End dialogue{]}} \\
\\ 
\texttt{At the end of this dialogue, Lorelei would feel...} \\
\texttt{Dismay} \\
\texttt{Exasperation} \\
\texttt{Amusement} \\
\texttt{Acceptance} \\
\\ 
\texttt{Give each of these possible emotions a score from 0{-}10 for the relative intensity that they are likely to be feeling each.} \\
\\ 
\texttt{You must output in the following format, including headings (of course, you should give your own scores), with no additional commentary:} \\
\\ 
\texttt{Dismay: <score>} \\
\texttt{Exasperation: <score>} \\
\texttt{Amusement: <score>} \\
\texttt{Acceptance: <score>} \\
\\ 
\\ 
\texttt{{[}End of answer{]}} \\
\\ 
\texttt{Remember: zero is a valid score, meaning they are likely not feeling that emotion. You must score at least one emotion > 0.} \\
\\ 
\texttt{Your answer:} \\
            }
        }
    }
     \\ \\

    \theutterance \stepcounter{utterance}  
    & \multicolumn{4}{p{0.6\linewidth}}{
        \cellcolor[rgb]{0.561,0.737,0.561}{
            \makecell[{{p{\linewidth}}}]{
                \texttt{\tiny{[P1$\rangle$GM]}}
                \texttt{Dismay: 2} \\
\texttt{Exasperation: 9} \\
\texttt{Amusement: 0} \\
\texttt{Acceptance: 1} \\
\\ 
\texttt{Lorelei's tone throughout the dialogue is increasingly frustrated and annoyed as Harold repeatedly dismisses her concerns about the future and the necessity of action. She uses phrases like "That's not the point," "You think throwing money...?" and "And I'm saying..." which indicate a growing irritation with Harold's perspective. While she remains firm in her logic (not amusement) and doesn't give up on her cause entirely (not dismay or acceptance), the primary emotional driver is the frustration of having her arguments consistently rejected by someone she views as out of touch.} \\
            }
        }
    }
    & & \\ \\

    \theutterance \stepcounter{utterance}  
    & & & \multicolumn{2}{p{0.3\linewidth}}{
        \cellcolor[rgb]{0.9,0.9,0.9}{
            \makecell[{{p{\linewidth}}}]{
                \texttt{\tiny{[GM$|$GM]}}
                \texttt{Dismay: 2} \\
\texttt{Exasperation: 9} \\
\texttt{Amusement: 0} \\
\texttt{Acceptance: 1} \\
            }
        }
    }
    & & \\ \\

    \theutterance \stepcounter{utterance}  
    & & & \multicolumn{2}{p{0.3\linewidth}}{
        \cellcolor[rgb]{0.9,0.9,0.9}{
            \makecell[{{p{\linewidth}}}]{
                \texttt{\tiny{[GM$|$GM]}}
                \texttt{Dismay: 4} \\
\texttt{Exasperation: 8} \\
\texttt{Amusement: 0} \\
\texttt{Acceptance: 0} \\
            }
        }
    }
    & & \\ \\

    \theutterance \stepcounter{utterance}  
    & & & \multicolumn{2}{p{0.3\linewidth}}{
        \cellcolor[rgb]{0.9,0.9,0.9}{
            \makecell[{{p{\linewidth}}}]{
                \texttt{\tiny{[GM$|$GM]}}
                \texttt{game\_result = LOSE} \\
            }
        }
    }
    & & \\ \\

\end{supertabular}
}

\end{document}
